\begin{dang}{Tính toán sức căng bề mặt của chất lỏng}
	\setcounter{paragraph}{0}
\Anlythuyet{\paragraph{Lực căng bề mặt}
\begin{align*} 
\boder{f=\sigma.\ell} 
\end{align*}
\begin{itemize}
\item Khi khung có đường kính vòng trong và vòng ngoài lần lượt là d và D: $f=\sigma \pi (D+d)$
\item Khi khung có đường kính vòng trong và vòng ngoài bằng nhau và bằng d: $f=\sigma .2\pi d$.
\item Khi khung rất mảnh có hình vuông cạnh a: $f=\sigma.8a$
\end{itemize}
\paragraph{Tính lực cần thiết để nâng vật ra khỏi chất lỏng}
\begin{itemize}
\item Để nâng được vật: $F_K\ge P+f$.
\item Lực tối thiểu: $F_K=P+f$.
\end{itemize}
\paragraph{Bài toán về hiện tượng nhỏ giọt chất lỏng}
\begin{itemize}
\item Đầu tiên giọt chất lỏng to dần nhưng chưa rơi xuống.
\item Đúng lúc giọt chất lỏng rơi
\begin{align*}
P=f\Leftrightarrow mg=\sigma.\ell\Leftrightarrow V_0Dg=\sigma \pi d\Leftrightarrow \dfrac{V}{n}.D.g=\sigma \pi d
\end{align*}
\item Trong đó: 
\begin{itemize}
\item n là số giọt chất lỏng 
\item $V_0$ là thể tích của 1 giọt chất lỏng 
\item V là thể tích chất lỏng trong ống 
\item D là khối lượng riêng của chất lỏng 
\item $\ell=\pi d$ là chu vi miệng ống
\item d là đường kính ống 
\end{itemize}
\end{itemize}}
\end{dang}
